% Einheit 2 von 4 – Abstand Punkt–Ebene (bank/abstaende/e2.jsonl, zusammenbau v0.1)
%% TODO Zeitmarke Einheit 2: Marken-Zeile nicht lesbar
%% TODO Zweigzeile Teil 1: was hier gelernt wird (Satz in Schülersprache aus dem Einheitstitel) – Einheit 2
\einheitenkopf[e2]{Einheit 2 von 4 · Abstand Punkt–Ebene}
\zweigzeile{Abitur GK}
\setcounter{aufgabe}{11}

%% TODO Titel: Ich-kann-Satz fehlt in der Bank (Platzhalter: Kettenname) – Nr. 12
%% TODO Anweisung: ein Satz über der ersten Teilaufgabe fehlt in der Bank – Nr. 12
\begin{aufgabe}{Abstand Punkt–Ebene}
%% TODO \rechenplatz{n}: Schrittzahl des Grundfalls fehlt in der Bank (nur bei fester Schreibform, 2.3 b)
\begin{teile}
\teil Hin oder zurück? Wie weit ist P(2 | −1 | 3) von der Ebene E: $2x - 2y + z = 1$ entfernt? Kreuze an, rechne nicht. \\ \kreuz{Abstand berechnen} \\ \kreuz{Punkt zum vorgegebenen Abstand bestimmen}
\teil Abstand von P(4 | 0 | 3) zur Ebene E: $2x - y + 2z = 5$? d = \leerfeld
\teil Ein Rauchmelder sitzt im Punkt R(2 | 1 | 1) eines Raums, eine Dachschräge liegt in der Ebene D: $2x + y + 4z = 21$; 1 LE = 1 m. Gib eine Gleichung der Geraden k an, auf der R und der Punkt der Dachschräge mit dem kleinsten Abstand zu R liegen. \hfill (Abitur 2017 LK) \\ k: $\vec x = $ \leerfeld
\teil Die Gerade g: $\vec x = (1 | 0 | 2) + t \cdot (2 | -2 | 1)$ steht senkrecht auf der Ebene E: $2x - 2y + z = 4$, der Punkt (1 | 0 | 2) liegt in E. Gib einen Punkt auf g an, der von E den Abstand 6 hat. \hfill (Abitur 2026 LK) \\ ( \leerfeld | \leerfeld | \leerfeld )
\teil Eine Kamera gleitet an einer Stange, die im Modell die Strecke RG mit R(2 | 3 | 0) und G(2 | 3 | 6) ist; 1 LE = 1 m. Sie soll 2 m von einer Glaswand in der Ebene E: $x + 2y + 2z = 20$ entfernt sein. Wo auf der Stange sitzt sie? \hfill (Abitur 2018 LK) \\ X( \leerfeld | \leerfeld | \leerfeld )
\teil Gegeben ist die Ebenenschar $E_k$: $kx + 2y - z - k = 0$. Vorgelegt ist die Lösung einer Aufgabe: $\frac{|k \cdot 1 + 2 \cdot 3 - 1 - k|}{\sqrt{k^2 + 5}} = 1 \Leftrightarrow k = -2\sqrt{5} \vee k = 2\sqrt{5}$. Formuliere eine passende Aufgabenstellung. \hfill (Abitur 2018 LK)
\end{teile}
\end{aufgabe}

%% TODO Titel: Ich-kann-Satz fehlt in der Bank (Platzhalter: Kettenname) – Nr. 13
%% TODO Anweisung: ein Satz über der ersten Teilaufgabe fehlt in der Bank – Nr. 13
\begin{aufgabe}{Abstand Punkt–Ebene – Fehler finden, Begründen, Anwendung}
\begin{teile}
\teil Finde den Fehler. Gesucht ist der Abstand von P(1 | 2 | 4) zur Ebene E: $2x + y + 2z = 3$. \rechnung{d &= |2 + 2 + 8 - 3| = 9}
\teil Warum liefert das Teilen durch den Betrag des Normalenvektors den Abstand? Begründe.
\teil Eine Drohne schwebt in P(4 | 6 | 9) über einer Dachfläche in der Ebene D: $3x + 4z = 40$; 1 LE = 1 m. Vorgeschrieben sind mindestens 2 m Abstand zum Dach. Hält die Drohne ihn ein? \janein
\end{teile}
\end{aufgabe}
